% Options for packages loaded elsewhere
\PassOptionsToPackage{unicode}{hyperref}
\PassOptionsToPackage{hyphens}{url}
\documentclass[
  mongolian,
]{article}
\usepackage{xcolor}
\usepackage{amsmath,amssymb}
\setcounter{secnumdepth}{-\maxdimen} % remove section numbering
\usepackage{iftex}
\ifPDFTeX
  \usepackage[T1]{fontenc}
  \usepackage[utf8]{inputenc}
  \usepackage{textcomp} % provide euro and other symbols
\else % if luatex or xetex
  \usepackage{unicode-math} % this also loads fontspec
  \defaultfontfeatures{Scale=MatchLowercase}
  \defaultfontfeatures[\rmfamily]{Ligatures=TeX,Scale=1}
\fi
\usepackage{lmodern}
\ifPDFTeX\else
  % xetex/luatex font selection
\fi
% Use upquote if available, for straight quotes in verbatim environments
\IfFileExists{upquote.sty}{\usepackage{upquote}}{}
\IfFileExists{microtype.sty}{% use microtype if available
  \usepackage[]{microtype}
  \UseMicrotypeSet[protrusion]{basicmath} % disable protrusion for tt fonts
}{}
\makeatletter
\@ifundefined{KOMAClassName}{% if non-KOMA class
  \IfFileExists{parskip.sty}{%
    \usepackage{parskip}
  }{% else
    \setlength{\parindent}{0pt}
    \setlength{\parskip}{6pt plus 2pt minus 1pt}}
}{% if KOMA class
  \KOMAoptions{parskip=half}}
\makeatother
\usepackage{color}
\usepackage{fancyvrb}
\newcommand{\VerbBar}{|}
\newcommand{\VERB}{\Verb[commandchars=\\\{\}]}
\DefineVerbatimEnvironment{Highlighting}{Verbatim}{commandchars=\\\{\}}
% Add ',fontsize=\small' for more characters per line
\newenvironment{Shaded}{}{}
\newcommand{\AlertTok}[1]{\textcolor[rgb]{1.00,0.00,0.00}{\textbf{#1}}}
\newcommand{\AnnotationTok}[1]{\textcolor[rgb]{0.38,0.63,0.69}{\textbf{\textit{#1}}}}
\newcommand{\AttributeTok}[1]{\textcolor[rgb]{0.49,0.56,0.16}{#1}}
\newcommand{\BaseNTok}[1]{\textcolor[rgb]{0.25,0.63,0.44}{#1}}
\newcommand{\BuiltInTok}[1]{\textcolor[rgb]{0.00,0.50,0.00}{#1}}
\newcommand{\CharTok}[1]{\textcolor[rgb]{0.25,0.44,0.63}{#1}}
\newcommand{\CommentTok}[1]{\textcolor[rgb]{0.38,0.63,0.69}{\textit{#1}}}
\newcommand{\CommentVarTok}[1]{\textcolor[rgb]{0.38,0.63,0.69}{\textbf{\textit{#1}}}}
\newcommand{\ConstantTok}[1]{\textcolor[rgb]{0.53,0.00,0.00}{#1}}
\newcommand{\ControlFlowTok}[1]{\textcolor[rgb]{0.00,0.44,0.13}{\textbf{#1}}}
\newcommand{\DataTypeTok}[1]{\textcolor[rgb]{0.56,0.13,0.00}{#1}}
\newcommand{\DecValTok}[1]{\textcolor[rgb]{0.25,0.63,0.44}{#1}}
\newcommand{\DocumentationTok}[1]{\textcolor[rgb]{0.73,0.13,0.13}{\textit{#1}}}
\newcommand{\ErrorTok}[1]{\textcolor[rgb]{1.00,0.00,0.00}{\textbf{#1}}}
\newcommand{\ExtensionTok}[1]{#1}
\newcommand{\FloatTok}[1]{\textcolor[rgb]{0.25,0.63,0.44}{#1}}
\newcommand{\FunctionTok}[1]{\textcolor[rgb]{0.02,0.16,0.49}{#1}}
\newcommand{\ImportTok}[1]{\textcolor[rgb]{0.00,0.50,0.00}{\textbf{#1}}}
\newcommand{\InformationTok}[1]{\textcolor[rgb]{0.38,0.63,0.69}{\textbf{\textit{#1}}}}
\newcommand{\KeywordTok}[1]{\textcolor[rgb]{0.00,0.44,0.13}{\textbf{#1}}}
\newcommand{\NormalTok}[1]{#1}
\newcommand{\OperatorTok}[1]{\textcolor[rgb]{0.40,0.40,0.40}{#1}}
\newcommand{\OtherTok}[1]{\textcolor[rgb]{0.00,0.44,0.13}{#1}}
\newcommand{\PreprocessorTok}[1]{\textcolor[rgb]{0.74,0.48,0.00}{#1}}
\newcommand{\RegionMarkerTok}[1]{#1}
\newcommand{\SpecialCharTok}[1]{\textcolor[rgb]{0.25,0.44,0.63}{#1}}
\newcommand{\SpecialStringTok}[1]{\textcolor[rgb]{0.73,0.40,0.53}{#1}}
\newcommand{\StringTok}[1]{\textcolor[rgb]{0.25,0.44,0.63}{#1}}
\newcommand{\VariableTok}[1]{\textcolor[rgb]{0.10,0.09,0.49}{#1}}
\newcommand{\VerbatimStringTok}[1]{\textcolor[rgb]{0.25,0.44,0.63}{#1}}
\newcommand{\WarningTok}[1]{\textcolor[rgb]{0.38,0.63,0.69}{\textbf{\textit{#1}}}}
\usepackage{longtable,booktabs,array}
\usepackage{calc} % for calculating minipage widths
% Correct order of tables after \paragraph or \subparagraph
\usepackage{etoolbox}
\makeatletter
\patchcmd\longtable{\par}{\if@noskipsec\mbox{}\fi\par}{}{}
\makeatother
% Allow footnotes in longtable head/foot
\IfFileExists{footnotehyper.sty}{\usepackage{footnotehyper}}{\usepackage{footnote}}
\makesavenoteenv{longtable}
\ifLuaTeX
\usepackage[bidi=basic,shorthands=off,,provide=*]{babel}
\else
\usepackage[bidi=default,shorthands=off,,provide=*]{babel}
\fi
\ifLuaTeX
  \usepackage{selnolig} % disable illegal ligatures
\fi
\setlength{\emergencystretch}{3em} % prevent overfull lines
\providecommand{\tightlist}{%
  \setlength{\itemsep}{0pt}\setlength{\parskip}{0pt}}
\usepackage{fontspec}
\setmainfont{Liberation Serif}
\setsansfont{Liberation Sans}
\setmonofont{Liberation Mono}[Scale=0.86]
\usepackage[a4paper,left=3cm,right=2cm,top=2cm,bottom=2cm,includefoot,headheight=15pt]{geometry}
\usepackage{ragged2e,indentfirst,xurl}
\usepackage{longtable,booktabs,array,calc,graphicx}
\usepackage{xcolor,fancyhdr,titlesec,enumitem,float}
\definecolor{ink}{HTML}{19354A}
\definecolor{accent}{HTML}{167D8D}
\titleformat{\section}{\Large\sffamily\bfseries\color{ink}}{}{0pt}{}
\titleformat{\subsection}{\large\sffamily\bfseries\color{ink}}{}{0pt}{}
\titleformat{\subsubsection}{\fontsize{11.5}{13.5}\selectfont\sffamily\bfseries\color{ink}}{}{0pt}{}
\titlespacing*{\section}{1.25cm}{0pt}{8pt}
\titlespacing*{\subsection}{1.25cm}{8pt}{4pt}
\titlespacing*{\subsubsection}{1.25cm}{6pt}{3pt}
\setlength{\parindent}{1.25cm}
\setlength{\parskip}{3pt}
\setlength{\emergencystretch}{3em}
\setlength{\tabcolsep}{4pt}
\renewcommand{\arraystretch}{1.08}
\setlist{nosep,leftmargin=1.5em,topsep=3pt}
\pagestyle{fancy}\fancyhf{}
\fancyhead[L]{\small\sffamily ICSI438 / SEMINAR 06}
\fancyhead[R]{\small\sffamily Nogoolin / MN}
\fancyfoot[R]{\small\thepage}
\renewcommand{\headrulewidth}{0.3pt}
\AtBeginDocument{\fontsize{10.5}{12.5}\selectfont\justifying\setlength{\parindent}{1.25cm}}


\AtBeginDocument{\renewcommand{\figurename}{Зураг}}
\usepackage{bookmark}
\IfFileExists{xurl.sty}{\usepackage{xurl}}{} % add URL line breaks if available
\urlstyle{same}
\hypersetup{
  pdflang={mn},
  hidelinks,
  pdfcreator={LaTeX via pandoc}}

\author{}
\date{}

\begin{document}

\hypersetup{
  pdftitle={Програм хангамжийн төслийн баримт бичиг — Seminar 6 Code Documentation},
  pdfauthor={Ц. Тэнгис}
}
\pagenumbering{gobble}
\thispagestyle{empty}
\noindent
\begin{minipage}[t][\dimexpr\textheight-\topskip\relax][t]{\linewidth}
  \begin{center}
    \textbf{МОНГОЛ УЛСЫН ИХ СУРГУУЛЬ}\\
    \textbf{МЭДЭЭЛЛИЙН ТЕХНОЛОГИ, ЭЛЕКТРОНИКИЙН СУРГУУЛЬ}\\
    \textbf{МЭДЭЭЛЭЛ, КОМПЬЮТЕРЫН УХААНЫ ТЭНХИМ}\\[1.7cm]
    \includegraphics[width=6cm]{../assets/num-logo.jpg}\\[1.2cm]
    {\fontsize{16}{20}\selectfont\bfseries
      ПРОГРАМ ХАНГАМЖИЙН ТӨСЛИЙН БАРИМТ БИЧИГ\par}
    \vspace{0.45cm}
    {\large\bfseries SEMINAR 6 — CODE DOCUMENTATION БА AI AUDIT}
  \end{center}

  \vspace{1.2cm}
  \begin{tabular}{>{\raggedright\arraybackslash}m{5cm}@{\hspace{1cm}}>{\raggedright\arraybackslash}m{\dimexpr\linewidth-6cm-2\tabcolsep\relax}}
    Хичээл: & ICSI438 — Software Project Documentation \\[0.55cm]
    Үндсэн төсөл: & Nogoolin \\[0.55cm]
    Хичээлийн багш: & Ph.D. Batnyam Battulga \\[0.55cm]
    Гүйцэтгэсэн: & Програм хангамжийн IV ангийн оюутан\\
                  & Ц. Тэнгис (22B1NUM6249) \\[0.55cm]
    Баримтын хувилбар: & v0.1 draft
  \end{tabular}

  \vfill
  \begin{center}
    Улаанбаатар\\[0.25cm]
    2026 он
  \end{center}
\end{minipage}
\clearpage
\pagenumbering{arabic}
\setcounter{page}{1}


\section{1. Кодын API
documentation}\label{ux43aux43eux434ux44bux43d-api-documentation}

Nogoolin-ийн backend service-үүдийг JSDoc-оор тайлбарлаж, TypeDoc-оор
HTML reference үүсгэв. Үндсэн төсөл TypeScript тул handout-ийн
JSDoc/TypeDoc хувилбарыг сонгосон. UE-6-ийн Python class-ийн оронд
inquiry service-ийн factory болон method-уудыг ашиглав.

Sem6-ийн эх ба файлууд:
\url{https://github.com/Tengis01/Nogoolin/tree/main/docs/ICSI438/sem6}.
\href{reference/index.html}{reference/} дотор generated API reference,
\href{audit/provenance.md}{audit/} дотор харьцуулалт,
\href{tools/package.json}{tools/} дотор build/test хэрэгслүүд бий.

\subsection{1.1 Хамрах хүрээ ба өмнөх ажлын
холбоо}\label{ux445ux430ux43cux440ux430ux445-ux445ux4afux440ux44dux44d-ux431ux430-ux4e9ux43cux43dux4e9ux445-ux430ux436ux43bux44bux43d-ux445ux43eux43bux431ux43eux43e}

Cart, category, inquiry, media, product, settings гэсэн зургаан service
модулийг хамруулсан. Нийт 44 symbol: 6 factory, 23 method, 6 service
type, 2 media interface, 7 interface field. Эдгээрийн \textbf{44/44 нь
тайлбартай}; function/method бүр parameter, return, error, жишээтэй.
Controller, database adapter, hook болон frontend нь энэ reference-ийн
entry point биш. Төслийн бүх export-ийн бүрэн coverage гэж үзэхгүй.

Sem3-ийн inquiry шаардлага нь service-ийн submit үйлдэлтэй, Sem4-ийн
Controller → Service → Repository бүтэц нь тайлбарын хариуцлагын
заагтай, Sem5-ийн inquiry endpoint нь controller-оор энэ service дуудах
урсгалтай холбогдоно. HTTP contract ба code reference хоёрын зорилго
ялгаатай.

\subsection{1.2 Inquiry service-ийн нэг
жишээ}\label{inquiry-service-ux438ux439ux43d-ux43dux44dux433-ux436ux438ux448ux44dux44d}

Жишээнүүдэд Ногоон Дарь эх бүтээгдэхүүнтэй холбоотой inquiry ашиглав.
Service factory-д repository dependency өгнө. UUID болон
repository-уудыг caller урьдчилан бэлтгэсэн байх нөхцөлийг reference-ийн
эхэнд тайлбарласан.

\begin{Shaded}
\begin{Highlighting}[]
\CommentTok{/**}
\CommentTok{ * Inquiry үүсгэнэ; product\_id байвал бүтээгдэхүүнийг шалгана.}
\CommentTok{ * }\AnnotationTok{@param}\CommentTok{ }\CommentVarTok{input}\CommentTok{ {-} Controller{-}оор schema шалгасан input.}
\CommentTok{ * }\AnnotationTok{@param}\CommentTok{ }\CommentVarTok{customerId}\CommentTok{ {-} JWT subject UUID эсвэл зочинд null.}
\CommentTok{ * }\AnnotationTok{@returns}\CommentTok{ Repository{-}ийн үүсгэсэн Inquiry.}
\CommentTok{ * }\AnnotationTok{@throws}\CommentTok{ }\CommentVarTok{PRODUCT\_NOT\_FOUND}\CommentTok{ (404), repository failure.}
\CommentTok{ */}
\KeywordTok{const}\NormalTok{ inquiry }\OperatorTok{=} \ControlFlowTok{await}\NormalTok{ service}\OperatorTok{.}\FunctionTok{submit}\NormalTok{(\{}
\NormalTok{  customer\_name}\OperatorTok{:} \StringTok{"Тест хэрэглэгч"}\OperatorTok{,}
\NormalTok{  phone}\OperatorTok{:} \StringTok{"99112233"}\OperatorTok{,}
\NormalTok{  message}\OperatorTok{:} \StringTok{"Ногоон Дарь эх байгаа юу?"}
\NormalTok{\}}\OperatorTok{,} \KeywordTok{null}\NormalTok{)}\OperatorTok{;}
\end{Highlighting}
\end{Shaded}

Input validation болон authentication нь controller/schema-д хийгддэг.
Service нь JWT шалгахгүй. Non-exported helper-үүдийг implementation
detail гэж тайлбарлаж, шууд public API болгон гаргаагүй. Зургаан
service-ийн comment засварын өмнөх ба дараах emitted JavaScript ижил
гарсан тул runtime үйлдэл өөрчлөгдөөгүй.

\clearpage

\section{2. AI audit ба эхтэй тулгасан
засвар}\label{ai-audit-ux431ux430-ux44dux445ux442ux44dux439-ux442ux443ux43bux433ux430ux441ux430ux43d-ux437ux430ux441ux432ux430ux440}

Bhatti Ch.5-ийн Explained, Concise, Clear, Usable, Trustworthy гэсэн
таван зарчмыг inquiry module-д хэрэглэв. Нэг controlled faulty draft-ийг
эхтэй тулгасан baseline-тай харьцуулж diff, audit хүснэгт,
Use--Verify--Cite log хадгалсан.

\subsection{2.1 Bhatti-ийн таван
зарчим}\label{bhatti-ux438ux439ux43d-ux442ux430ux432ux430ux43d-ux437ux430ux440ux447ux438ux43c}

Доорх 0--5 оноо нь локал self-assessment; багшийн болон peer-ийн үнэлгээ
биш.

\begin{longtable}[]{@{}
  >{\raggedright\arraybackslash}p{(\linewidth - 4\tabcolsep) * \real{0.3333}}
  >{\raggedright\arraybackslash}p{(\linewidth - 4\tabcolsep) * \real{0.3333}}
  >{\raggedright\arraybackslash}p{(\linewidth - 4\tabcolsep) * \real{0.3333}}@{}}
\toprule\noalign{}
\begin{minipage}[b]{\linewidth}\raggedright
Зарчим
\end{minipage} & \begin{minipage}[b]{\linewidth}\raggedright
Draft → baseline
\end{minipage} & \begin{minipage}[b]{\linewidth}\raggedright
Засварын гол үр дүн
\end{minipage} \\
\midrule\noalign{}
\endhead
\bottomrule\noalign{}
\endlastfoot
Explained & 2 → 4 & Validation/auth болон service behavior-ийн зааг
тодорхой. \\
Concise & 4 → 4 & Method бүрийн зорилгыг богино бичсэн. \\
Clear & 2 → 4 & Array ба pagination wrapper-ийг ялгасан. \\
Usable & 2 → 4 & Domain жишээ, dependency нөхцөл, snippet type-check. \\
Trustworthy & 1 → 4 & Source ба local test-тэй тулгасан; live test
хийгдээгүй. \\
\end{longtable}

Дундаж нь draft 2.2/5, baseline 4.0/5. Бүрэн audit нь тусдаа Markdown
файлд байна.

\subsection{2.2 Гурван
засвар}\label{ux433ux443ux440ux432ux430ux43d-ux437ux430ux441ux432ux430ux440}

\textbf{C1 --- Validation.} Draft нь submit утсыг яг найман цифр эсэхийг
өөрөө шалгана гэсэн. Бодит implementation-д ийм шалгалт байхгүй.
Controller input schema хэрэглэж, service зөвхөн product байгаа эсэхийг
шалган create дууддаг гэж зассан. Sem5-д тэмдэглэсэн SRS--phone schema
drift хэвээр байгаа; энэ ажлаар бүтээгдэхүүний дүрэм өөрчлөөгүй.

\textbf{C2 --- Return.} Draft нь listOwn pagination metadata буцаана
гэсэн. Бодит return нь Inquiry array; хоосон үед хоосон array буцаадгийг
test шалгасан.

\textbf{C3 --- Status ба transaction.} Draft нь зөвхөн new → contacted →
closed дараалалтай, update ба audit нэг transaction гэсэн. Service нь
шилжилтийн дарааллыг хориглодоггүй. Audit алдаа гарсан ч өмнөх status
update rollback болохгүйг local double test шалгасан. Анхны кодын
lifecycle comment-ийг мөн бодит implementation-тэй нийцүүлсэн.

\subsection{2.3 Гарал ба hallucination
test}\label{ux433ux430ux440ux430ux43b-ux431ux430-hallucination-test}

Baseline нь AI-assisted, эх кодтой тулгасан хувилбар. Controlled
draft-ийн алдааг зориуд audit exercise-д оруулсан. Үүнийг Тэнгис бие
даан гараар бичсэн эсвэл GPT-4/Claude-ийн санамсаргүй алдаа олсон гэж
үзэхгүй. Одоогийн assistant нь Codex/GPT-6.

Reflection test-д байхгүй getInquiryById method болон InquiryReceipt
type-ийг negative fixture болгон өгсөн. Test эдгээрийг илрүүлж,
corrected reference list-ийг зөвшөөрсөн. Энэ нь detector ажиллах
нотолгоо; бодит model output-оос invented API олсон шалгуурын нотолгоо
биш. Independent manual review болон тухайн output-ийн audit нээлттэй.

\clearpage

\section{3. Чанар ба автомат
build}\label{ux447ux430ux43dux430ux440-ux431ux430-ux430ux432ux442ux43eux43cux430ux442-build}

\subsection{3.1 Chinchilla-ийн зөвлөгөөг хэрэглэсэн
нь}\label{chinchilla-ux438ux439ux43d-ux437ux4e9ux432ux43bux4e9ux433ux4e9ux4e9ux433-ux445ux44dux440ux44dux433ux43bux44dux441ux44dux43d-ux43dux44c}

Handout ``Six Principles'' гэж нэрлэсэн ч зургаан нэрийг жагсаагаагүй.
Доорх checklist нь Ch.3 болон Ch.7-ийн зөвлөгөөг энэ ажилд хэрэглэхээр
нэгтгэсэн.

\begin{longtable}[]{@{}
  >{\raggedright\arraybackslash}p{(\linewidth - 2\tabcolsep) * \real{0.5000}}
  >{\raggedright\arraybackslash}p{(\linewidth - 2\tabcolsep) * \real{0.5000}}@{}}
\toprule\noalign{}
\begin{minipage}[b]{\linewidth}\raggedright
Шалгалт
\end{minipage} & \begin{minipage}[b]{\linewidth}\raggedright
Энэ ажилд хэрэглэсэн нь
\end{minipage} \\
\midrule\noalign{}
\endhead
\bottomrule\noalign{}
\endlastfoot
Нэг мөр нэр томьёо ба domain & Ногоон Дарь эх → inquiry; тогтмол
method/type нэр. \\
Үйлдэл, хариуцагч тодорхой & Controller шалгалт ба service үйлдлийг
ялгасан. \\
Богино, хэрэгтэй тайлбар & Summary, parameter, result, error
дараалалтай. \\
Бодит хэрэглээтэй жишээ & Guest inquiry болон admin status update. \\
Prerequisite тодорхой & Dependency, UUID, Buffer-ийн нөхцөлтэй. \\
Жишээг шалгаж хадгалах & Snippet type-check, local test, strict
build. \\
\end{longtable}

\subsection{3.2 TypeDoc ба GitLab
CI}\label{typedoc-ux431ux430-gitlab-ci}

TypeDoc 0.28.14, TypeScript 5.7.3 болон холбогдох dependency-г Sem6-ийн
lockfile-д тогтоосон. Build нь compiler error, missing documentation
болон invalid link-ийг шалгана; warning гарвал амжилтгүй дуусна.
Scope-оос гаднах repository/schema contract-д зориулан notExported
шалгалтыг унтраасан бөгөөд шалтгааныг reference index-д бичсэн.

\begin{Shaded}
\begin{Highlighting}[]
\NormalTok{npm ci {-}{-}prefix docs/ICSI438/sem6/tools}
\NormalTok{npm run verify {-}{-}prefix docs/ICSI438/sem6/tools}
\NormalTok{python3 docs/ICSI438/sem6/latex/build.py}
\end{Highlighting}
\end{Shaded}

Командуудыг Nogoolin root-оос ажиллуулна. Verifier нь coverage → test →
TypeDoc гэсэн дарааллаар ажиллаж verification log хадгална. TeX-ийн
compile нь Markdown-оос дахин үүсгэхгүй; Ctrl+S нь LaTeX Workshop-оор
PDF шинэчилнэ.

GitLab pipeline-д docs-build job HTML reference үүсгэж public artifact
гаргана. Default branch дээр publish-docs job тэр artifact-ийг GitLab
Pages-д нийтлэхээр тохируулсан. Root CI файл нь Sem6-ийн CI файлыг
include хийдэг. Тусдаа GitLab repo ашиглахад шаардлагатай
source/config-ийг сонгон export хийх хэрэгсэл мөн бэлдсэн.

\subsection{3.3 Бодит
verification}\label{ux431ux43eux434ux438ux442-verification}

\begin{longtable}[]{@{}
  >{\raggedright\arraybackslash}p{(\linewidth - 2\tabcolsep) * \real{0.5000}}
  >{\raggedright\arraybackslash}p{(\linewidth - 2\tabcolsep) * \real{0.5000}}@{}}
\toprule\noalign{}
\begin{minipage}[b]{\linewidth}\raggedright
Шалгалт
\end{minipage} & \begin{minipage}[b]{\linewidth}\raggedright
Үр дүн
\end{minipage} \\
\midrule\noalign{}
\endhead
\bottomrule\noalign{}
\endlastfoot
Public symbol inventory & Service scope-д 44/44. \\
Зан төлөв ба snippet/reflection & 9 test давсан; snippet-ууд type-check
давсан. \\
TypeDoc HTML & Local build амжилттай, warning байхгүй. \\
Runtime өөрчлөгдсөн эсэх & Зургаан service-ийн JavaScript output
ижил. \\
CI YAML & Parse, include path, job dependency бүтэц шалгасан. \\
GitLab pipeline / Pages & Project байхгүй тул ажиллуулаагүй. \\
\end{longtable}

\clearpage

\section{4. Technical debt ба
дүгнэлт}\label{technical-debt-ux431ux430-ux434ux4afux433ux43dux44dux43bux442}

\subsection{4.1 TODO/FIXME
бүртгэл}\label{todofixme-ux431ux4afux440ux442ux433ux44dux43b}

Backend, packages болон apps-ийн TS/TSX эхээс нэг TODO олдсон: mobile
auth session AsyncStorage-д хадгалагдаж байгаа бөгөөд encrypted
storage-д шилжүүлэх хэрэгтэй. Хуучин Phase 2 хугацааны тэмдэглэгээг авч,
local issue draft-тай холбосон. Authentication кодын үйлдлийг
өөрчлөөгүй.

Issue draft-д session encryption, restore/refresh, logout cleanup,
existing session migration, storage failure болон Android/iOS
verification-ийн acceptance criteria бий. GitLab project байхгүй учраас
issue ID зохиогоогүй; бодит issue үүсэх хүртэл энэ шалгуур нээлттэй.
Файл нь issues хавтас дотор байна.

\subsection{4.2 Reflection}\label{reflection}

\textbf{Bhatti-аас авсан санаа.} Миний хувьд Trustworthy зарчим хамгийн
хэрэгтэй байлаа. Зөв хэлбэртэй JSDoc өөрөө зөв тайлбарын баталгаа биш.
Жишээлбэл status-ийн сумтай тайлбарыг хараад шилжилтийн хориг байна гэж
ойлгож болох ч код нь түүнийг хэрэгжүүлээгүй. Тайлбарын claim бүрийг
implementation эсвэл test-тэй холбох нь ийм зөрүүг илрүүлнэ.

\textbf{Chinchilla-аас үргэлжлүүлэх зүйл.} Нэг domain example-ийг баримт
даяар үргэлжлүүлэх нь method бүрт шинэ нөхцөл ойлгох ачааллыг багасгаж
байна. Цаашид жишээнүүдийн dependency нөхцөлийг эхэнд нь тайлбарлаж,
type-check болон test-ийг build-тэй холбоно.

\textbf{Ном ба бодит ажлын ялгаа.} Номонд жишээг шалгаж, review хийж,
нийтлэн хадгалах workflow ярьдаг. Энэ ажилд локал build/test бэлэн
болсон боловч GitLab project, бие даасан manual review болон бодит AI
output-ийн харьцуулалтын нотолгоо бүрэн биш. Эдгээрийг локал PDF-ээр
орлуулахгүй гэж үзлээ.

\subsection{4.3 Үлдсэн
шалгуур}\label{ux4afux43bux434ux441ux44dux43d-ux448ux430ux43bux433ux443ux443ux440}

GitLab CI green, live Pages URL, бодит issue дугаар болон
repository-wide public symbol coverage батлагдаагүй. UE-6-ийн
independent manual authoring, GPT-4/Claude response болон spontaneous
invented API audit мөн нээлттэй. Одоогийн багц нь TypeScript хувилбарын
локал implementation ба controlled audit exercise юм.

\subsection{4.4 Эх
сурвалж}\label{ux44dux445-ux441ux443ux440ux432ux430ux43bux436}

\begin{itemize}
\tightlist
\item
  Sprint 06 Lab Instructions, х.1--2: US-6.1--6.4, UE-6, DoD.
\item
  Bhatti et al., Docs for Developers: Ch.4 х.79; Ch.5 х.86--94.
\item
  Chinchilla, Technical Writing for Software Developers: Ch.3 х.28--33;
  Ch.6 х.69--82; Ch.7 х.89--92; Ch.9 х.123.
\item
  \href{https://typedoc.org/documents/Options.Validation.html}{TypeDoc
  validation}.
\item
  \href{https://docs.gitlab.com/user/project/pages/getting_started/pages_from_scratch/}{GitLab
  Pages workflow}.
\item
  Nogoolin source, Sem6 symbol inventory, audit diff ба verification
  log.
\end{itemize}

\end{document}
